\section{Introducción}

En esta unidad, trataremos con una técnica básica de modelación predictiva llamada \emph{regresión lineal}, la cuál permite crear un modelo a partir de una base de datos histórica.

Nuestro propósito es entender las matemáticas detrás de la regresión lineal e ilustrar sus resultado a través de su implementación en varias bases de datos.

\subsection{Contexto e importancia}

La regresión lineal es una de las técnicas más fundamentales y ampliamente utilizadas en estadística y ciencia de datos. Desarrollada en el siglo XIX por Francis Galton y formalizada matemáticamente por Karl Pearson y Ronald Fisher, la regresión lineal continúa siendo una herramienta esencial en:

\begin{itemize}
	\item \textbf{Predicción:} Estimar valores futuros basándose en datos históricos (ventas, precios, demanda).
	\item \textbf{Inferencia:} Comprender las relaciones entre variables y cuantificar su impacto.
	\item \textbf{Control de procesos:} Optimizar sistemas industriales y de negocio.
	\item \textbf{Investigación científica:} Probar hipótesis sobre relaciones causales.
\end{itemize}

A pesar de su simplicidad, la regresión lineal sirve como base para técnicas más avanzadas como regresión polinomial, modelos lineales generalizados, y métodos de aprendizaje automático.

\subsection{Mapa de ruta}
\begin{itemize}
	\item Las matemáticas detrás de la regresión lineal.
	\item Implementación de la regresión lineal con Python.
	\item Interpretación de los parámetros resultantes.
	\item Validación del modelo.
	\item Manejo de problemas relacionados con regresión lineal.
\end{itemize}


\subsection{Modelos matemáticos}
Un \emph{modelo matemático/estadístico/predictivo} es una ecuación matemática que consiste en \emph{entradas} que producen \emph{salidas} cuando el valor de las variables entrantes se introduce en el modelo.

Los modelos pueden clasificarse en:
\begin{itemize}
	\item \textbf{Determinísticos:} Producen siempre el mismo resultado para las mismas entradas (ej: $F = ma$).
	\item \textbf{Estocásticos:} Incorporan incertidumbre y variabilidad aleatoria (ej: modelos de regresión).
\end{itemize}

La regresión lineal es un modelo estocástico porque reconoce que las observaciones reales no siguen perfectamente una relación lineal, sino que incluyen un componente de error aleatorio.

\subsection{Ejemplo}
Por ejemplo, supongamos que el precio $P$ de una casa es \emph{linealmente dependiente} en su tamaño $S$, comodidades $A$ y disponibilidad de transporte $T$.



La ecuación correspondiente sería
\begin{align}
	P = a_{1}\times S+ a_{2} \times A + a_{3} \times T + \epsilon
\end{align}

donde $\epsilon$ representa el error aleatorio (factores no considerados en el modelo).


Esta ecuación es llamado el \emph{modelo} y los coeficientes $a_{1},a_{2},a_{3}$ son sus parámetros.



La variable $P$ es resultado predicho, mientras que $S,A,T$ con las variables de entrada, que son datos conocidos.


Sin embargo, los parámetros $a_{i}$ deben ser estimados a partir de los datos históricos.


Una vez que estos parámetros son determinados, el modelo está listo para ser evaluado.

\subsection{Relación con la correlación}

Es importante distinguir entre \emph{correlación} y \emph{regresión}:

\begin{itemize}
	\item \textbf{Correlación:} Mide la fuerza y dirección de la relación lineal entre dos variables. Es simétrica ($corr(X,Y) = corr(Y,X)$) y adimensional (varía entre -1 y 1).
	
	\item \textbf{Regresión:} Modela la relación funcional entre variables, permitiendo predecir valores de una variable a partir de otras. No es simétrica y los coeficientes tienen unidades.
\end{itemize}

Una alta correlación entre variables sugiere que una regresión lineal podría ser apropiada, pero la correlación por sí sola no implica causalidad ni proporciona un modelo predictivo.

\subsection{Supuestos del modelo}

Para que un modelo de regresión lineal sea válido y sus inferencias confiables, deben cumplirse ciertos supuestos:

\begin{enumerate}
	\item \textbf{Linealidad:} La relación entre las variables predictoras y la respuesta es lineal.
	\item \textbf{Independencia:} Las observaciones son independientes entre sí.
	\item \textbf{Homocedasticidad:} La varianza de los errores es constante.
	\item \textbf{Normalidad:} Los errores siguen una distribución normal (especialmente importante para muestras pequeñas).
	\item \textbf{No multicolinealidad:} Los predictores no están altamente correlacionados entre sí.
\end{enumerate}

En secciones posteriores discutiremos cómo verificar estos supuestos y qué hacer cuando no se cumplen.

\subsection{Tipos de regresión lineal}

Dependiendo del número de variables predictoras, distinguimos:

\begin{itemize}
	\item \textbf{Regresión lineal simple:} Una sola variable predictora ($Y = \beta_0 + \beta_1 X + \epsilon$).
	\item \textbf{Regresión lineal múltiple:} Dos o más variables predictoras ($Y = \beta_0 + \beta_1 X_1 + \beta_2 X_2 + \cdots + \beta_p X_p + \epsilon$).
\end{itemize}

Ambos tipos utilizan el mismo principio fundamental: el método de mínimos cuadrados para estimar los parámetros que minimizan la suma de los errores cuadráticos.

\subsection{Aplicaciones prácticas}

Algunos ejemplos de aplicaciones de regresión lineal incluyen:

\begin{itemize}
	\item \textbf{Finanzas:} Predecir precios de acciones, evaluar riesgo crediticio, estimar retornos de inversión.
	\item \textbf{Marketing:} Modelar el impacto de la publicidad en ventas, analizar efectividad de campañas.
	\item \textbf{Medicina:} Relacionar dosis de medicamentos con respuestas, predecir progresión de enfermedades.
	\item \textbf{Ingeniería:} Control de calidad, predicción de fallas, optimización de procesos.
	\item \textbf{Ciencias sociales:} Estudiar relaciones entre variables socioeconómicas, educativas, políticas.
\end{itemize}

En las siguientes secciones desarrollaremos la teoría matemática y mostraremos cómo implementar y validar modelos de regresión lineal usando Python.

